\Subhead{5.2.1}{Геометрическое и дифракционное приближения. Влияние неровностей земной поверхности. Земные радиоволны}

На трассе с размерами много больше $\lambda$ поле удобно представить суммой прямого, отражённого и преломлённых лучей. Но близко к границе тени и за препятствием существенна дифракция. Для препятствия между передатчиком и приёмником масштаб первой зоны Френеля
\begin{equation}
 r_F=\sqrt{\frac{\lambda d_1d_2}{d_1+d_2}}.
\end{equation}
Он получается из превышения оптического пути примерно на $\lambda/2$. Даже при геометрической видимости препятствие в значимой части зоны меняет амплитуду. Дифракция на кромке даёт переход в тень через интегралы Френеля; полная тень геометрической оптики не означает точно нулевого поля.

Отражение от земли зависит от комплексной проницаемости, проводимости и поляризации. Для гладкой поверхности применяются коэффициенты Френеля; прямой и отражённый лучи складываются по фазам, а не по мощностям. При гауссовской высоте $\sigma_h$ зеркальная амплитуда ослабляется из-за случайных фаз $2kh\cos\theta$. Земная кривизна и рельеф меняют видимость и дифракционные пути; плоская двухлучевая модель имеет ограниченную область применимости.

Поверхностная земная волна связана с границей воздух–грунт и огибает Землю с потерями. Она особенно важна на длинных волнах; хорошая проводимость морской поверхности уменьшает потери. Вертикальная поляризация обычно выгоднее для дальнего поверхностного распространения. Это не просто один отражённый луч, а граничная волновая задача.
